%!TEX root = newmain.tex

  The technique of Sections~\ref{sec:translation}
  and~\ref{sec:dispatchtable} extends to cyclic contracts with
  separate cycles such as \stipula{Deposit}. Here we examine the
  additional challenges posed by cycles and clarify the resulting loss
  of precision.

  In \stipula{Deposit}, the \stipula{Farm} and the \stipula{Client}
  may repeatedly invoke \stipula{send} and \stipula{buy}, forming a
  cycle in the contract underlying automaton
  (Figure~\ref{fig:deposit-fsm}). The cycle is exited through a
  dedicated event executed at most once. This reflects a common
  contractual pattern: (\emph{i}) cycles have a unique entry/exit
  state and contain only function invocations; (\emph{ii}) cycles are
  not nested; (\emph{iii}) events occur outside cycles and typically
  enforce their termination. As a consequence, each event is executed
  at most once, even in the presence of cyclic behaviour.
  
  This observation allows us to reuse the dispatch table mechanism
  developed for acyclic contracts, since the set of events remains
  finite.  However, events generated inside cycles require special
  care: although they share the same index, different iterations may
  assign them different time values.  To handle this, we assign such
  events an \emph{abstract} clock value, encoded as \jml{-2}.  Rather
  than reasoning about their precise scheduling time, we only record
  that they have been generated. We illustrate the resulting loss
  of precision with an example. Consider the cyclic \Stipula\ function
  declaration:
  
  {\small
\begin{lstlisting}[numbers=none,frame=none,mathescape=true,language=stipula,escapechar=§]
@Q f()[] {
   now + 1 $\event$ @Q {} $\tostate$ @Q1
   now + 2 $\event$ @Q {} $\tostate$ @Q2
} $\tostate$ @Q
\end{lstlisting}
  }

  According to the \Stipula\ semantics, the second event is never
  triggered, because the first one always precedes it, so state
  \stipula{@Q2} is unreachable. But since the events occur in a cycle,
  both are given the abstract clock value \jml{-2}, enabling an
  infeasible execution in the \JML/\Java\ model that triggers the
  second event and reaches \stipula{@Q2}.

  For modeling cycles in \JML\ we refine $\updateDT{\Q}{W}$ and
  distinguish whether the generating function of an event lies inside
  or outside a cycle. If it belongs to a cycle, the events in $W$ are
  inserted into the dispatch table with the abstract time value
  \jml{-2}, instead of a concrete timestamp.  Accordingly, the
  predicate \ensureDT$_{\f}$ is adapted to enforce that \emph{exactly}
  the events generated by $\f$ are assigned the value \jml{-2}, while
  all other entries of \jml{DT} remain unchanged. The structure of the
  specification remains the same as in the acyclic case, differing
  only in the assigned time value. The remaining definitions in
  Section~\ref{sec:dispatchtable} are omitted, because they are
  straightforward.

  The crucial adaptation concerns the extension of
  $\genscenario{\C}{\Q}$ (Listing~\ref{tab:fungenerate}) to states
  within cycles. For such states, we use the notation
  $\genscenarioplus{\C}{\Q,{\tt count}}{\Q'}$, where $\Q$ denotes the
  entry/exit state of the cycle and \jml{count} is the iteration
  counter introduced by the translation (see
  Listing~\ref{tab:fungeneratecycles}).  Each cycle is translated into
  a dedicated {\Java} \jml{while} loop whose body encodes all possible
  traversals of the cycle. A counter tracks the number of iterations,
  bounded by a cycle-specific constant \stipula{MAX_ITE}$_\Q$. This
  bound becomes a parameter of the \stipula{behaviour} method in
  Listing~\ref{tab:translate} with precondition %
  \jml{requires MAX_ITE$_\Q\geq\,$0} and provides a natural
  abstraction: although {\Stipula} allows unbounded cyclic behaviour
  (\emph{e.g.}, in \stipula{Deposit}), only finitely many interactions
  occur in practice.  Finally, the time-progress variable
  \stipula{u_Q} used in the acyclic case is omitted for cyclic states,
  since events generated within cycles carry an undetermined time
  value and do not require explicit clock advancement.

{\small
\begin{lstlisting}[frame=none,numbers=none,language={[JML]Java},mathescape=true,caption={The function $\genscenarioplus{\C}{\Q,{\tt count}}{\Q'}$.},label={tab:fungeneratecycles},float,escapechar=�]
�{\textrm{There are two cases:\\
  \hspace*{-2em}(1) $\Q$ is the entry/exit state of a cycle; ${\f}_1, \ldots, {\f}_n$ and\ 
  $\mathtt{ev_{\jmath\,1}}, \ldots, {\tt ev_{\jmath\,m}}$ are functions\\
  \hspace*{-2em} with initial state $\Q$ and $\texttt{final}$ states\ 
  $\Q_1, \ldots, \Q_n$ and $\Q_1', \ldots, \Q_m'$, respectively.}�

$\genscenarioplus{\C}{}{\Q}  \; \eqdef$
   int entry = minTimeEntry({ev$_{\jmath 1}$, $\cdots$ , ev$_{\jmath m}$}) ; 
   if ((entry == -1) || (DT[entry] > now && w_Q > 0)) {
     int count = 0 ;
     /*@ loop_invariant 0 <= count <= MAX_ITE$_\Q$ ; // bounds
       @ loop_invariant ... ; // assets, fields, and DT invariants     
       @ decreases MAX_ITE$_\Q$ - count ; // termination expression
       @*/
     while (count < MAX_ITE$_\Q$) {
       if (w_Q == 1) { $\f_1(\vect{\tt x_1}, \vect{\tt k_1})$ ; if (Q != Q$_1$) $\genscenarioplus{\C}{\Q,{\tt count}}{\Q_1}$ }
       else if (w_Q == 2) { $\f_2(\vect{\tt x_2}, \vect{\tt k_2})$ ; if (Q != Q$_2$) $\genscenarioplus{\C}{\Q,{\tt count}}{\Q_2}$ }
       ...
       else { $\f_n(\vect{{\tt x}_n}, \vect{{\tt k}_n})$ ;  if (Q != Q$_n$) $\genscenarioplus{\C}{\Q,{\tt count}}{\Q_n}$ }  // max_Q == n
       count = count + 1 ;
     } 
     entry = minTimeEntry({ev$_{\jmath 1}$,$\cdots$,ev$_{\jmath m}$}) ; // we are back in state Q
     if (DT[entry] > now || DT[entry] == -2) { 
       if (DT[entry] > now) now = DT[entry] ;
       if (entry == ev$_{\jmath 1}$) { event_${\jmath}$1() ; $\genscenarioplus{\C}{}{\Q_1'}$ }
       else if (entry == ev$_{\jmath 2}$) { event_${\jmath}$2(); $\genscenarioplus{\C}{}{\Q_2'}$ }
       ...
       else { event_${\jmath}$m() ; $\genscenarioplus{\C}{}{\Q_m'}$ }      //  entry == ev$_{\jmath m}$
     } 
   } else { if (DT[entry] > now) now = DT[entry] ;
     if (entry == ev$_{\jmath 1}$) { event_${\jmath}$1() ; $\genscenarioplus{\C}{}{\Q_1'}$ }
     else if (entry == ev$_{\jmath 2}$) { event_${\jmath}$2(); $\genscenarioplus{\C}{}{\Q_2'}$ }
     ...
     else { event_${\jmath}$m() ; $\genscenarioplus{\C}{}{\Q_m'}$ }      //  entry == ev$_{\jmath m}$
   }    
     
 �\textrm{(2) ${\Q}'$ is a state of a cycle that is \emph{not} the entry/exit state. Let ${\f}_1, \ldots, {\f}_\ell$ be the functions whose final states $\Q_1, \ldots, \Q_\ell$ are in the cycle.}�

$\genscenarioplus{\C}{\Q,{\tt count}}{\Q'} \eqdef$ if (w_Q$'$ == 1) { $\f_1(\vect{\tt x_1},\vect{\tt k_1})$ ; if (Q != Q$_1$) $\genscenarioplus{\C}{\Q,{\tt count}}{\Q_1}$ }
                else if (w_Q$'$ == 2) { $\f_2(\vect{\tt x_2}, \vect{\tt k_2})$ ; if (Q != Q$_2$) $\genscenarioplus{\C}{\Q,{\tt count}}{\Q_2}$ }
                ...
                else { $\f_\ell(\vect{\tt x_\ell}, \vect{\tt k_\ell})$ ; if (Q != Q$_\ell$) $\genscenarioplus{\C}{\Q,{\tt count}}{\Q_\ell}$ } 
\end{lstlisting}
}

From a deductive verification perspective,
% using counters to define loop iterations has an important advantage:
% it simplifies the generation of loop invariants in {\KeY}.  In
% practice, the invariant of \jml{while} loop can relate the current
% state of the contract, the dispatch table, and the iteration
% counter, thus reducing the burden of invariant discovery and
% enabling {\KeY} to reason effectively about cyclic behaviors within
% a bounded yet sound approximation.  Of course,
finding suitable loop invariants is a very difficult task in formal
program verification and usually not amenable to automation. For this
reason our translation technique leaves the supply of post
condition-specific invariants to the verification engineer.
%
For example, in the case of \stipula{Deposit}, using the following
invariant ($v_s$ is the actual value of the parameters {\tt h} and
{\tt k} in {\tt begin} and {\tt send}; $v_b$ is the value of {\tt w}
in {\tt buy}) it is possible to derive in {\KeY} the property that all
the flour acquired by the \stipula{Client} has been paid to the
\stipula{Farm} at agreed cost.

{\small
\begin{lstlisting}[frame=none,numbers=none,language={[JML]Java},mathescape=true]
loop_invariant (Client_flour == \old(Client_flour) && flour == \old(flour)+$v_{s}$)$\,$||
               (Client_flour == \old(Client_flour)+$v_b$/cost_flour &&
                Farm == \old(Farm)+$v_b$) ;
\end{lstlisting}   }   

  Another limitation of our current encoding is that function
  parameters are treated as symbolic constants throughout a loop,
  meaning all iterations are analysed under a single symbolic
  instantiation (e.g., $v_s$ and $v_b$). While {\Stipula} allows
  parameters to vary between iterations, verifying this
  non-determinism would require loop invariants over sequences of
  iteration-dependent values, which tools like {\KeY} cannot currently
  handle automatically. Supporting such cases is theoretically
  possible with user-supplied invariants and interactive proofs, but
  will further reduce the automation of our approach.

For the cyclic case, we can prove a soundness theorem for the contracts analysed here, 
assuming constraints on repeated function calls with identical arguments and on 
the number of iterations.

\begin{theorem}[Soundness]
  \label{thm.soundnessacyclic}
  Let $\C$ be a {\Stipula} contract
  % $\contract {\comma} 0$ be the initial configuration of a
  % {\Stipula} contract {\C}
  with separate cycles and events that are outside cycles. Let
  \[
    \contract {\comma} {\tt 0} \lred{(\vect A,\, \vect{A_i}{:} \vect{u_i}{\,}^{i \in 1..\ell})}
    \lred{}^* \contract_1 {\comma} \Time_1 \lred{\mu_1} \lred{}^* 
    \contract_2 {\comma} \Time_2 \lred{\mu_2} \lred{}^* \cdots 
    % \contract_{n} {\comma} \Time_n
    \lred{\mu_{n+1}} \lred{}^* \contract_{n+1} {\comma} \Time_{n+1}
  \]
  where $\Q_i$ are the states of $\contract_i$, with $i \in 1..n+1$,
  and $\mu_i$ are either function invocations
  $A_i{:}\f_i(\vect{v_i})[\vect{v_i'}]$ or events
  $\ev_i$. Additionally, if $\mu_i$ and $\mu_j$ are invocations of the
  same function then $\mu_i = \mu_j$ (the invocations have the same
  actual values).
  % 
  Then, %with suitable values of the constants \stipula{MAX_ITE}$_\Q$
  the above computation may be traced in the code of
  $\genscenario{\C}{\cdot}$ as follows. Consider
  $\contract_i {\comma} \Time_i \lred{\mu_i} \lred{}^* \contract_{i+1}
  {\comma} \Time_{i+1}$ and assume that, in the code of
  $\genscenario{\C}{\Q_i}$ or of
  $\genscenarioplus{\C}{\Q,{\tt count}}{\Q_i}$ (when $\Q_i$ is part of
  a cycle),
  \begin{enumerate}
  \item \stipula{now} = $\Time_i$ if $\mu_i$ is an event not generated
    inside a cycle; \stipula{now} $\leq \Time_i$, otherwise;
  \item if the multiset of pending events in $\contract_i$ is
    \[
      {\Time}_{\jmath_1} \,\event_{\ev_{\jmath_1}}\,
      \Qwithat_{\jmath_1}\, \{ \, S_{\jmath_1} \,\} \,\tostate\,
      \Qwithat_{\jmath_1}' ~|~ \cdots ~|~ {\Time}_{\jmath_\kappa}
      \event_{\ev_{\jmath_k}}\, \Qwithat_{\jmath_k}\, \{ \,
      S_{\jmath_k} \,\} \,\tostate\, \Qwithat_{\jmath_k}'
    \]
    and the dispatch table is such that
    ${\tt DT}[\ev_{\jmath_i}] \in \{ {\Time}_{\jmath_i}, \mbox{\tt
      -2}\}$, with $i \in 1.. k$.
  \end{enumerate}
  Then, if $\mu_i = A_i{:}\f_i(\vect{v_i})[\vect{v_i'}]$ it is
  possible to execute the method $\f_i(\vect{v_i},\vect{v_i'})$ and if
  $\mu_i = \ev_{\jmath_i}$, it is possible to execute method
  ${\tt event}\_{\jmath_i}(~)$ and obtain a state where (1.) and (2.)
  hold for $\contract_{i+1} \comma \Time_{i+1}$.
\end{theorem}

This result crucially depends on the values of \stipula{MAX_ITE}$_\Q$.
If reliable upper bounds on the number of possible iterations are
known, the analysis carried out with {\KeY} is guaranteed to explore
all the legally admissible behaviour of the contract. However, in
contrast to the acyclic setting, completeness can no longer be
claimed.  As discussed at the beginning of the section, the introduction of over-approximations, specifically,
assigning the abstract time value {\tt -2} to events generated inside
cycles, may lead the translated {\Java} program to exhibit behaviour
not admissible by the semantics of {\Stipula}. As a consequence, the
verification performed by {\KeY} may encounter spurious executions,
potentially preventing the proof of properties that do hold in the
original contract.

%%% Local Variables:
%%% mode: latex
%%% TeX-master: "newmain"
%%% End:
